\section{Background and Related Work}
\label{sec:pipeline}

\subsection{Windowed reconstruction and its three fixed budgets}

The published method \cite{wqpe2025,modularshor2025} replaces one
$\ell$-qubit phase register with $W$ overlapping windows of width
$\ell_w\ll\ell$, each measured independently by its own block-local QPE
circuit, and reassembles the phase from the per-window measurement counts in
three classical steps: generate the top-$k$ candidate bit patterns per
window, stitch compatible candidates across overlaps under a carry-aware
consistency check with a beam of width $P$, and recover the order by
continued fractions. Factor extraction follows from an even order $r$ via
$\gcd(a^{r/2}\pm1,N)$.

Three quantities in that stage --- the retained candidate count $k$, the beam
width $P$, and the per-window shot count $S_w$ --- are fixed at configuration
time and identical for every window of every instance. They are the subject
of both halves of this article: Part I establishes that they are the binding
constraint, and Part II evaluates replacing them.

\subsection{Problem formulation: the reproduction, and what counts as success}

All results here are from an independent from-scratch reproduction, not from
the original authors' code. Throughout, a trial \emph{succeeds} only if it
recovers a correct non-trivial factor pair of $N$ end to end. This is a
deliberately unforgiving criterion: it does not credit a partially correct
phase, and it makes every reported rate directly comparable across the two
halves.

One consequence of it shapes Part I. At the small multiplicative orders our
compilation ceiling permits (Sec.~\ref{sec:cliff}), the continued-fraction
multiplier search is extremely forgiving --- theoretically under-resolved
configurations still recover the correct order. Any characterisation that
sweeps one parameter at a time will therefore find nothing, and
Sec.~\ref{sec:ofat} shows that it does.

% =====================================================================
